Air pollution is a severe public health problem in Delhi, where the official Air Quality Index (AQI) exceeded 300 (\emph{Very Poor} or worse) on 22.7\% of days between 2015 and 2026. This project forecasts Delhi's next-day AQI using a Random Forest regressor optimized by Particle Swarm Optimization (PSO), a nature-inspired metaheuristic based on the social behaviour of bird flocks. The dataset combines 4,170 days of official Central Pollution Control Board (CPCB) AQI bulletins (May 2015 to October 2026) with daily ERA5 weather reanalysis, giving 32 leakage-free features of AQI history, meteorology and season. PSO was implemented from first principles with linearly decreasing inertia, velocity clamping and reflecting boundaries. Each particle encodes six Random Forest hyperparameters and 32 threshold-based feature switches, so the swarm tunes the model and selects its inputs at the same time, scored by walk-forward validation error. Training used 2015 to 2023, and the 1,026 days from January 2024 to October 2026 were held out for testing. The joint search reduced validation RMSE by 4.3\% (41.52 to 39.73), twice the gain of a hyperparameter-only search, kept 14 of 32 features, and outperformed random search with the same 312-evaluation budget. On the test period, the selected model reached an RMSE of 33.77, an MAE of 25.15 and an $R^2$ of 0.882, improving on a naive persistence forecast by 16.0\%; averaged over five runs, joint PSO gave a test RMSE of 33.47, 1.0\% below the default Random Forest. The selected model's gain over the default forest was small (0.25\%) and not statistically significant ($p = 0.76$), and it was slightly worse than the hyperparameter-only model (RMSE 33.57), revealing that the larger search partly overfitted its validation folds. An untuned ridge regression performed comparably (RMSE 33.45), showing that informative features matter more than tuning for this problem. The work provides a reproducible, tested framework and an honest assessment of what nature-inspired optimization contributes to air quality forecasting.
